% Looks like lecture numbering has been bumped.
% This lecture should come after lecture 12 (22 September), but before lecture 14, which is currently labelled lecture13.
% It's confusing, but the section numbering will actually work out!
\section{Lecture 13: Maps to Projective Space (again)}

\begin{theorem}[Arbitrary Morphisms to {$\P^n$}]\label{arb:morphism_to_Pn}
    Maps from an arbitrary scheme $X$ to $\P^n$ are equivalent to two pieces of data:
    \begin{enumerate}
        \item A line bundle $\mathcal{L}$ on $X$; and,
        \item A \emph{locally split injection} $L\xhookrightarrow{}\mathcal{O}_X^{n+1}$
    \end{enumerate}
\end{theorem}

The line bundle is locally trivial on some open cover $\{U_i\}$, which is that $\mathcal{L}\mid_{U_i}\cong\mathcal{O}_X\mid_{U_i}$, and we refer to the choice of isomorphism as a \emph{trivialisation}.

Consider an $\mathcal{O}_x$-linear map $\mathcal{O}_X\to\mathcal{O}_X^{n+1}$.
This map must be $\mathcal{O}_X$-linear (that is, $\phi_U(as)=a\phi_U(s)$ for all $a\in\mathcal{O}_X(U)$).
On this basis, once we know where $1$ goes, we know where everything goes by linearity:
\begin{figure}[H]
    \centering
    \begin{tikzcd}[row sep=0]
        \mathcal{O}_X\ar[r,"\varphi"]&\mathcal{O}_X^{n+1}\\
        1\ar[r, maps to]&(f_0,\dots,f_n)\\
        g\ar[r, maps to]&(gf_0,\dots,gf_n)
    \end{tikzcd}
\end{figure}

If it is \emph{split} there is a map `backwards' on our injection.
How do we make this precise?

In the case above, consider $\psi:\mathcal{O}_X\xleftarrow{}\mathcal{O}_X^{n+1}$.
A standard basis for this map is
\begin{align*}
    e_0&\coloneqq(1,0,\dots,0)\\
    e_1&\coloneqq(0,1,\dots,0)\\
    \dots\\
    e_{n}&\coloneqq(0,\dots,1)
\end{align*}
and these give a map by sending $e_i\mapsto g_i$ for some $g_i$ in $\mathcal{O}_X$.
Clearly, the image of the $e_i$s in $\mathcal{O}_X$ determines what happens on the maps of sheaves.

\begin{definition}[Splitting]
    Without reference to Exact Sequences, we make the definition that a \emph{split map} is one such that the following map is an identity on $A$, that is, $\psi\circ\varphi=\operatorname{id}_A$.
    \begin{figure}[H]
        \centering
        \begin{tikzcd}
            A\ar[r,shift left, "\varphi"]& B\ar[l, shift left, "\psi"]
        \end{tikzcd}
    \end{figure}

    Note that $\varphi$ surjective (and, being an identity on the round trip) implies the existence of an isomorphism, but we do not make this assumption and only operate on the image of $\varphi$ in $B$.
\end{definition}

\begin{proposition}
    For a map $\varphi$, $\psi$ is a splitting of $\varphi$ if and only if $1=f_0g_0+\dots+f_ng_n$.
    \begin{proof}
        We start by noting that the map is $\mathcal{O}_X$-linear, and take $\varphi(a)=(af_0,\dots,af_n)$.
        Then, note that the map in the reverse direction is $\psi(a_0,\dots,a_n)=\sum a_ig_i$.
        It is clear that $(\psi\circ\varphi)(a)=a\sum f_ig_i$.
        We get the identity here if and only if $\sum f_ig_i=1$.
    \end{proof}
\end{proposition}

\begin{example}
    Consider the map $\varphi:\Z\to\Z\oplus\Z$, defined by $1\mapsto(2,3)$.
    We can consider this as a map of $\mathcal{O}_X$-modules $\mathcal{O}_{\Spec{\Z}}(\Spec{\Z})\to\mathcal{O}_{\Spec{\Z}}^{\oplus 2}(\Spec{\Z})$.
    At the point $(7)\in\Spec{\Z}$ the stalk $\mathcal{O}_{\Spec\Z,(7)}$ is the local ring $\Z_{(7)}$, which contains $\frac{1}{2}$.
    This gives the reverse map as $\psi(e_1)=\frac{1}{2}$, and $\psi(e_2)=0$.
    At the stalk $\mathcal{O}_{\Spec\Z,(2)}$, $\psi(e_1)=0$, and $\psi(e_2)=\frac{1}{3}$.
    
    Note that at every point of $\Spec{\Z}$, at least one of $2$ or $3$ is a unit.
\end{example}

\begin{definition}[Locally Split Injection]
    A map $\mathcal{L}\to\mathcal{O}_X^{n+1}$ is a \emph{locally split injection} if and only if, when evaluated at every point of $x\in X$:
    \begin{align*}
        f:\faktor{\mathcal{L}_x}{\frak{m}_x\mathcal{L}_x}\to\left(\faktor{\mathcal{O}_{X,x}}{\frak{m}_x}\right)^{n+1}=\kappa{(x)}^{n+1}
    \end{align*}
    is non-zero.
\end{definition}

The stalk $\mathcal{L}_x$ of the sheaf $\mathcal{L}$ is a free $\mathcal{O}_X$-module of rank one.
Evaluation at $x$ sends a section $s\in\mathcal{L}_x$ to its image in $\faktor{\mathcal{L}_x}{\frak{m}_x\mathcal{L}_x}$ (note the similarity to evaluation in $\mathcal{O}_x$).
We expect a one dimensional vector space.

\begin{nonexample}
    Consider the map 
    \begin{figure}[H]
        \centering
            \begin{tikzcd}[row sep=0]
                \Z\ar[r]&\Z\oplus\Z\\
                1\ar[r,maps to]&(4,6)
            \end{tikzcd}
    \end{figure}

    This is an injection (the kernel is trivial), but it is not an injection when evaluated at $(2)$:
    \begin{figure}[H]
        \centering
            \begin{tikzcd}[row sep=0]
                \faktor{\Z}{2\Z}\ar[r]&\faktor{\Z}{2\Z}\oplus\faktor{\Z}{2\Z}\\
                1\ar[r,maps to]&(0,0)
            \end{tikzcd}
    \end{figure}
\end{nonexample}

\begin{definition}[Locally Split Injection]
    An injection $\iota:A\xhookrightarrow{}B$ of $\mathcal{O}_X$-modules is locally split if for some open cover $\{U_i\}$ there exists $r$ such that $r\circ\iota=\operatorname{id}_A$ on each $U_i$.
\end{definition}
\begin{example}
    Any split injection is locally split.
\end{example}

Suppose we're given a scheme $X$ and a locally split injection for a line bundle $\mathcal{L}\xhookrightarrow{}\mathcal{O}_X^{n+1}$ and that we want to define a morphism of schemes $X\to\P^n$.
First, we can define the map locally:
Take some $x\in X$, choose a neighbourhood, and define the map to $\P^n$, and then we can build an open cover of $X$.
Then we must check for agreement on the overlaps.
Choose some $U$ containing $x$ such that $\mathcal{L}\mid_U$ is trivial.
Such a $U$ exists: by the definition of a line bundle, it is locally free of rank $1$.
To define the map globally, we will need to make choices and later check that this will work globally.

For our choice of $U$, choose a trivialisation $\mathcal{L}\mid_U\cong\mathcal{O}_X$.
This will allow us to write the map in local coordinates:
\begin{figure}[H]
    \centering
        \begin{tikzcd}[row sep=0]
            \mathcal{O}_U\ar[r, hook]&\mathcal{O}^{n+1}\\
            1\ar[r,maps to]&(f_0,\dots,f_n)
        \end{tikzcd}
\end{figure}

We know that at the stalk $\mathcal{O}_{X,x}$, at least one of the $f_i$ is a unit.
Without loss of generality, assume $f_0$ is a unit on all of $U$.
This gives us that $f_0$ is a regular function on $U$.

Then we may define a map $U\to\Spec{\Z[\frac{x_1}{x_0},\dots,\frac{x_n}{x_0}]}\subseteq\P^n$.
Note that $\Spec{\Z[\frac{x_1}{x_0},\dots,\frac{x_n}{x_0}]}$ is one of the standard charts for $\P^n$, and is open as it forms our previously discussed standard open cover of $\P^n$.

We now recall that to define a map $U\to\Spec{R}$ (for $R$ a ring), we just need to specify a map on rings:
\begin{figure}[H]
    \centering
        \begin{tikzcd}[row sep=0]
            \Z[\frac{x_1}{x_0},\dots,\frac{x_n}{x_0}]\ar[r]&\mathcal{O}_U(U)\\
            \frac{x_i}{x_0}\ar[r,maps to]&\frac{f_i}{f_0}
        \end{tikzcd}
\end{figure}

Does this map depend on the choices that we have made?

Choice 1: Which trivialisation to select? We know that isomorphisms between rings must be related by a unit:
Choose a second isomorphism and multiply by some unit $g$ on $U$.
Then $1\mapsto(f_0,\dots,f_n)$ becomes $1\mapsto(gf_0,\dots,gf_n)$.
This gives the same ring map: $\frac{gx_i}{gx_0}=\frac{x_i}{x_0}$.
We can think of this as `the slope is unchanged'.

Does the choice of unit matter here?
It now matters how the $\A_i$s were glued together.
We could have chosen a different $f_i$ that is a unit on $U$ (assuming one exists).
What happens if both $f_0$ and $f_1$ (without loss of generality) are units?
Consider the following maps (and associated subsets of $\P^n$):
\begin{figure}[H]
    \centering
        \begin{tikzcd}[row sep=0]
            &\Spec{\Z[\frac{x_1}{x_0},\dots,\frac{x_n}{x_0}]}\\
            U\ar[rd]\ar[ru]& &\P^n\ar[lu, phantom, sloped, "\subset"]\arrow[ld, phantom, sloped, "\subset"]\\
            &\Spec{\Z[\frac{x_0}{x_1},\dots,\frac{x_n}{x_1}]}
        \end{tikzcd}
\end{figure}

We can express the intersection of these two open affine subsets of $\P^n$ by thinking of them as a subring of a much bigger localisation.
Let:
\begin{align*}
    V_{01}&=\Spec{\Z[\frac{x_1}{x_0},\frac{x_0}{x_1},\frac{x_2}{x_0},\dots,\frac{x_n}{x_0}]}\\
    V_{10}&=\Spec{\Z[\frac{x_0}{x_1},\frac{x_1}{x_0},\frac{x_2}{x_0},\dots,\frac{x_n}{x_1}]}
\end{align*}

The map must land in both $V_{01}$ and $V_{10}$.
We identify these in the same way that we previously built $\P^n$.
For $U\to V_{01}$ (on the level of rings):
\begin{equation*}
    \frac{x_1}{x_0}\mapsto\frac{f_1}{f_0},\frac{x_0}{x_1}\mapsto\frac{f_0}{f_1},\frac{x_i}{x_0}\mapsto\frac{f_i}{f_0}
\end{equation*}
and for $U\to V_{10}$ (again, on the level of rings):
\begin{equation*}
    \frac{x_0}{x_1}\mapsto\frac{f_0}{f_1},\frac{x_1}{x_0}\mapsto\frac{f_1}{f_0},\frac{x_i}{x_1}\mapsto\frac{f_i}{f_1}
\end{equation*}

We note that we can send $\frac{x_2}{x_1}\mapsto\frac{x_2}{x_0}\frac{x_0}{x_1}=\frac{x_2}{x_1}$, which shows the relation by a unit (of the appropriate ring).
With the $\P^n$ identification of $V_{01}\cong V_{10}$, these two agree.
\begin{quote}
    Choosing the slope of the anti-slope doesn't matter because we have glued the slopes already.
\end{quote}

We now have a locally defined map to $\P^n$ on all open subsets:
We have maps $U_1\to\P^n$ and $U_2\to\P^n$ and at least two maps on $U_1\cap U_2\to\P^n$.
So $\mathcal{L}\mid_{U_1\cap U_2}\cong\mathcal{O}_{U_1\cap U_2}$ and we can choose a coordinate which is a unit.
On $U_1\cap U_2$, we can inherit this from either $U_1$ or $U_2$, and our choices on $U_1$ and $U_2$ don't need to match.
Our map to $\P^n$ does not depend on the choice of unit!

Different choices all lead to the same map on $\P^n$. Thus we have a map from all of $X$ to $\P^n$.
Note that for each $x\in X$, there may be multiple choices of $U$. The choice doesn't matter here, either.

Suppose we're given a map $X\to\P^n$. How do we find a line bundle and a locally split injection $\mathcal{L}\xhookrightarrow{}\mathcal{O}_X^{n+1}$?
We can construct this using the transition functions on the standard open cover.
Take $\P^n=\bigcup V_i$ with $V_i=\Spec{[\frac{x_0}{x_i},\dots,\frac{x_n}{x_i}]}$.
Let $U_i=\phi^{-1}(V_i)$, and this gives us an open cover of $X$ (recall that a map $X\to\P^n$ must be continuous on the underlying topology).
On $U_i$, we can take the trivial line bundle
\begin{figure}[H]
    \centering
    \begin{tikzcd}[row sep=0]
        \mathcal{O}_{U_i}\ar[r]&\mathcal{O}_{U_i}^{n+1}\\
        1\ar[r,maps to]&(h_0,\dots,h_{i-1},1,h_{i+1},\dots,h_n)
    \end{tikzcd}
\end{figure}

We note that $U_i\to V_i$ is a map of affine schemes, so we simply reverse the ring maps!
The topological pullback of $\frac{x_0}{x_i}$ is a function on $U_i$, and a map $U_i\to V_i$ determined by these functions.

Let $h_0,\dots,h_{i-1},h_{i+1},\dots,h_n$ in $\mathcal{O}_{U_i}$ be the images of $\frac{x_0}{x_i},\dots,\frac{x_n}{x_i}$.
Consider (in order to find a trivial line bundle on $\mathcal{O}_{U_i}$):
\begin{figure}[H]
    \centering
        \begin{tikzcd}[row sep=0]
            \mathcal{O}_{U_i}\ar[r]&\mathcal{O}_{U_i}^{n+1}\\
            1\ar[r,maps to]&(h_0,\dots,h_{i-1},1,h_{i+1},\dots,h_n)
        \end{tikzcd}
\end{figure}

This is a locally split injection.
On each $U_i$, we've now given a free rank $1$ sheaf and a locally split injection.
But note that we wanted a global object. What are $U_0$ and $U_1$ (etc)?
We need line bundles (sheaves that are locally free).

We have two pieces of data on $U_0$ and $U_1$ that obviously don't agree:
\begin{figure}[H]
    \centering
        \begin{tikzcd}[row sep=0]
            U_0\cap U_1:&\mathcal{O}_{U_0\cap U_1}\ar[r]&\mathcal{O}_{U_0\cap U_1}^{n+1}\\
            \text{from $U_0$:}&1\ar[r,maps to]&(1,h_1,h_2,\dots,h_n)\\
            \text{from $U_1$:}&1\ar[r,maps to]&(h_0^\prime,1,h_2^\prime,\dots,h_n^\prime)
        \end{tikzcd}
\end{figure}

There's no apriori reason for the two maps restricted from $U_0$ and $U_1$ to be equal.
Where did the $h_i$ and $h_i^\prime$ come from? $\frac{x_i}{x_0}\mapsto h_i$ and $\frac{x_i}{x_1}\mapsto h_i^\prime$.
As they are the same maps to projective space, they are related by relations on $\frac{x_i}{x_0}$ and $\frac{x_i}{x_1}$.
This gives us $h_0^\prime=h_1^{-1}$ (otherwise the map to projective space would not have been well-defined).

We immediately see that $\frac{x_2}{x_1}=\frac{x_2}{x_0}\frac{x_0}{x_1}\mapsto h_2^\prime$.
Then $h_i^\prime=h_2h_1^{-1}$.
If this had not been the case, the maps on the intersection would not have been well-defined.

Then we have:
\begin{figure}[H]
    \centering
        \begin{tikzcd}[row sep=0]
            \mathcal{O}_{U_0\cap U_1}\ar[r,"h"]&\mathcal{O}_{U_0\cap U_1}^{\oplus n+1}\\
            \mathcal{O}_{U_0\cap U_1}\ar[r,"h^\prime"]&\mathcal{O}_{U_0\cap U_1}^{\oplus n+1}
        \end{tikzcd}
\end{figure}

We don't glue these by the identity map as they are locally two different gadgets.
Consider:
\begin{figure}[H]
    \centering
        \begin{tikzcd}
            \mathcal{O}_{U_0\cap U_1}\ar[d]\ar[r,"h"]&\mathcal{O}_{U_0\cap U_1}^{\oplus n+1}\ar[d,sloped, phantom,"="]\\
            \mathcal{O}_{U_0\cap U_1}\ar[r,"h^\prime"]&\mathcal{O}_{U_0\cap U_1}^{\oplus n+1}
        \end{tikzcd}
\end{figure}

So on $U_0\cap U_1$, we get a transition function such that $h_{01}$ is the image of $\frac{x_1}{x_0}$.
We also gain a transition map on every double intersection $h_{ij}$ is the image of $\frac{x_j}{x_i}$, which is a unit on $U_i\cap U_j$, and these satisfy the cocycle condition $h_{ij}h_{jk}h_{ki}=1$.

Using these transition functions, we can define a locally free sheaf of rank 1, and this is over $\mathcal{L}$, and comes with a map $\mathcal{L}\xhookrightarrow{}\mathcal{O}_X^{\oplus n+1}$.
The detail is clearly messy, and we don't see it in any of our reference textbooks.

We now have a locally split injection split on the $U_i$s $\mathcal{L}\xhookrightarrow{}\mathcal{O}_X^{\oplus n+1}$ on $X$ and that's why we need the source to be a line bundle.

The global line may not be directly trivialisable, but we only need $X$ to be a locally ringed space (for evaluations): $X$ need not be a scheme.
We also used that $\P^n$ was a scheme.
